%%%%%%%%%%%%%%%%%%%%%%%%%%%%%%%%%%%%%%%%%%%%%%%%%%%%%%%%%%%%%%%%%%%%%%%%
% Part 3: Sections 6–8 (Experimental Setup & Stats, Results, Validation)
% (Poster Sections 6, 7, 8)
%%%%%%%%%%%%%%%%%%%%%%%%%%%%%%%%%%%%%%%%%%%%%%%%%%%%%%%%%%%%%%%%%%%%%%%%

\section{Experimental Setup and Statistical Methodology}
\label{sec:setup}

\subsection{Simulation Environment}

All experiments use the Heston stochastic volatility model with parameters in Table~\ref{tab:heston_params_setup}. Simulations run on a single NVIDIA GPU with PyTorch 2.0+ and CUDA 11.8+.

\subsection{Dataset Generation}

80,000 Heston paths generated via Euler--Maruyama full-truncation scheme: 50K training, 10K validation, 20K test. $N = 30$ daily time steps over a 30-day horizon. ATM European call option ($K = S_0 = 100$), proportional transaction cost $c_{\mathrm{tc}} = 0.001$.

\begin{table}[H]
\centering\small
\caption{Heston simulation parameters.}
\label{tab:heston_params_setup}
\begin{tabular}{@{}llr@{}}
\toprule
\textbf{Symbol} & \textbf{Description} & \textbf{Value} \\\midrule
$S_0 = K$ & Initial price / strike & 100 \\
$v_0$ & Initial variance & 0.04 \\
$\kappa$ & Mean-reversion speed & 1.0 \\
$\theta_v$ & Long-run variance & 0.04 \\
$\xi$ & Vol-of-vol & 0.2 \\
$\rho$ & Correlation & $-0.7$ \\
$r$ & Risk-free rate & 0.05 \\
$T$ & Maturity & 30/365 \\
$N$ & Time steps & 30 \\
$c_{\mathrm{tc}}$ & Transaction cost & 0.001 \\
\bottomrule
\end{tabular}
\end{table}

\subsection{Feature Engineering Pipeline}

At each time step $k$: $I_k = (S_k/S_0, \log(S_k/K), \sqrt{v_k}, \tau_k, \Delta_k^{\mathrm{BS}}) \in \R^5$, where $\tau_k = (T-t_k)/T$ and $\Delta_k^{\mathrm{BS}}$ is Black--Scholes delta at current implied volatility $\sqrt{v_k}$. RSE additionally computes regime features: $r_k = (\mathrm{RV}^{(5)}, \mathrm{RV}^{(10)}, \mathrm{RV}^{(20)}, \mathrm{Trend}, \mathrm{Mom}, \mathrm{Jump}) \in \R^6$.

\subsection{Training Protocol}

Identical two-stage (CVaR $\to$ Entropic) protocol for all models:
\begin{itemize}[noitemsep]
    \item \textbf{Stage 1:} 50 epochs, $\mathcal{L}_1 = \cvar_{0.95}(-\mathrm{P\&L})$, lr $= 10^{-3}$, patience 15
    \item \textbf{Stage 2:} 30 epochs, $\mathcal{L}_2 = \rho_\lambda(-\mathrm{P\&L}) + \gamma \sum_k |\Delta\delta_k|$, lr $= 10^{-4}$, patience 10
    \item Adam optimizer, weight decay $10^{-4}$, gradient clipping $\|\nabla\| \leq 5.0$, batch size 256
\end{itemize}

Novel model-specific additions operate \emph{within} this protocol: W-DRO-T adds DRO penalty only in Stage~2; 3SCH splits 80 epochs as 40+15+25; RSE pre-trains base models then trains gating.

\subsection{Hyperparameter Configuration}

\begin{table}[H]
\centering\scriptsize
\caption{Complete hyperparameter table across all models.}
\label{tab:hyperparams}
\begin{tabular}{@{}llccccc@{}}
\toprule
\textbf{Param} & \textbf{Desc.} & \textbf{LSTM} & \textbf{Trans.} & \textbf{W-DRO-T} & \textbf{3SCH} & \textbf{RSE} \\\midrule
Hidden dim & Architecture & 50 & 64 & 64 & 50 & 50/64 \\
Layers & Depth & 2 & 3 & 3 & 2 & 2/3 \\
Parameters & Count & 54K & 85K & 85K & 54K & $\sim$200K \\
$\delta_{\max}$ & Delta clamp & 1.5 & 1.5 & 1.5 & 1.5 & 1.5 \\
$\varepsilon$ & DRO radius & --- & --- & 0.1 & --- & --- \\
$K$ & Regimes & --- & --- & --- & --- & 4 \\
$\alpha_{\text{start}}/\alpha_{\text{end}}$ & Annealing & --- & --- & --- & 0.8/0.2 & --- \\
$\gamma$ & Trading penalty & $10^{-3}$ & $10^{-3}$ & $10^{-3}$ & $10^{-3}$ & $10^{-3}$ \\
Train time & Per seed & 273s & 7534s & 7119s & 295s & 924s \\
\bottomrule
\end{tabular}
\end{table}

\subsection{Statistical Methodology}
\label{sec:stats}

\subsubsection{Multi-Seed Experimental Protocol}

10 independent seeds: $s \in \{42, 142, 242, 342, 442, 542, 642, 742, 842, 942\}$. Each seed controls: data generation (Heston path sampling), weight initialisation (neural network parameters), batch ordering (data shuffling), and dropout masks (Transformer). This ensures that performance variation reflects genuine algorithmic robustness rather than lucky initialisations.

\subsubsection{Bootstrap Confidence Intervals}

For metric values $\hat\mu_1, \ldots, \hat\mu_{10}$ across seeds~\citep{efron1979bootstrap}:
\begin{enumerate}[noitemsep]
    \item Draw $B = 10{,}000$ bootstrap samples $(\hat\mu^*_{b,1}, \ldots, \hat\mu^*_{b,10})$ by resampling with replacement
    \item Compute $\bar\mu^*_b = \frac{1}{10}\sum_i \hat\mu^*_{b,i}$ for each resample
    \item The 95\% CI is $[\bar\mu^*_{(0.025B)},\; \bar\mu^*_{(0.975B)}]$
\end{enumerate}

\subsubsection{Paired Statistical Tests}

For model $A$ vs.\ baseline LSTM, let $d_s = \hat\mu_s^A - \hat\mu_s^{\mathrm{LSTM}}$ be the paired difference at seed $s$. We test $H_0: \E[d] = 0$ via paired $t$-test (or Wilcoxon signed-rank if Shapiro--Wilk rejects normality at $\alpha = 0.05$).

\subsubsection{Multiple Testing Correction}

Family-wise error control via Holm--Bonferroni~\citep{holm1979simple}: sort $p$-values $p_{(1)} \le \cdots \le p_{(K)}$ and reject $H_{0,(j)}$ if $p_{(j)} \le \alpha/(K - j + 1)$ where $\alpha = 0.05$.

\subsubsection{Effect Size Analysis}

Cohen's $d = \bar{d} / s_d$ where $\bar{d}$ is the mean paired difference and $s_d$ its standard deviation~\citep{cohen1988statistical}. Convention: $|d| < 0.2$ negligible, $0.2$--$0.5$ small, $0.5$--$0.8$ medium, $> 0.8$ large.

\section{Experimental Results}
\label{sec:results}

\subsection{Main Performance Metrics}

\begin{table}[H]
\centering\small
\caption{Test set performance (mean $\pm$ std, 10 seeds, 50K training paths). Lower is better for all metrics except trading volume.}
\label{tab:main_results}
\begin{tabular}{@{}lcccccr@{}}
\toprule
\textbf{Model} & \textbf{CVaR$_{95}$} & \textbf{Std P\&L} & \textbf{Entropic} & \textbf{Trade Vol} & \textbf{Mean P\&L} & \textbf{CV\%} \\\midrule
\textbf{RSE} & $\mathbf{3.109 \pm 0.010}$* & 0.456 & 2.376 & 1.379 & $-2.277$ & \textbf{0.30} \\
LSTM & $3.215 \pm 0.018$ & 0.449 & 2.379 & 1.137 & $-2.278$ & 0.55 \\
3SCH & $3.219 \pm 0.022$ & 0.453 & 2.381 & 1.122 & $-2.277$ & 0.69 \\
W-DRO-T & $3.227 \pm 0.024$ & 0.450 & 2.380 & 1.119 & $-2.277$ & 0.75 \\
Transformer & $3.234 \pm 0.030$ & 0.450 & 2.380 & 1.117 & $-2.277$ & 0.92 \\
\midrule

\end{tabular}
\\[3pt]{\footnotesize * $p < 0.0001$ vs LSTM after Holm--Bonferroni correction.}
\end{table}

RSE achieves the best CVaR$_{95}$ of $3.109 \pm 0.010$, a 3.29\% improvement over LSTM. Critically, RSE also achieves the lowest coefficient of variation (0.30\%), confirming that ensemble combination stabilizes performance across random seeds. All models achieve nearly identical mean P\&L ($\approx -2.277$), confirming negligible bias from ensemble averaging.

\subsection{Comparative Model Performance}

\begin{table}[H]
\centering\small
\caption{Paired comparison vs.\ LSTM (Holm--Bonferroni corrected).}
\label{tab:paired}
\begin{tabular}{@{}lrcccc@{}}
\toprule
\textbf{Model} & \textbf{$\Delta$\%} & \textbf{$p$-value} & \textbf{Cohen's $d$} & \textbf{Effect} & \textbf{Sig.} \\\midrule
\textbf{RSE} & $-3.29$ & $< 0.0001$ & $-7.41$ & Very Large & *** \\
3SCH & $+0.13$ & $0.438$ & $+0.20$ & Small & ns \\
W-DRO-T & $+0.37$ & $0.016$ & $+0.56$ & Medium & * \\
Transformer & $+0.59$ & $0.006$ & $+0.93$ & Large & ns \\
\bottomrule
\end{tabular}
\end{table}

RSE's improvement is highly significant against all baselines: $p < 0.0001$ with very large effect sizes (Cohen's $d$ from $-5.66$ to $-7.41$). 3SCH is statistically indistinguishable from LSTM ($p = 0.438$, $d = 0.20$). W-DRO-T's in-distribution CVaR is statistically indistinguishable from Transformer ($p = 0.188$, $d = -0.26$).

\subsection{Bootstrap Confidence Interval Analysis}

\begin{table}[H]
\centering\small
\caption{95\% bootstrap CIs (10,000 resamples).}
\label{tab:bootstrap}
\begin{tabular}{@{}lccc@{}}
\toprule
\textbf{Model} & \textbf{Mean} & \textbf{Lower} & \textbf{Upper} \\\midrule
RSE & 3.109 & 3.103 & 3.115 \\
LSTM & 3.215 & 3.204 & 3.226 \\
3SCH & 3.219 & 3.205 & 3.233 \\
W-DRO-T & 3.227 & 3.212 & 3.241 \\
Transformer & 3.234 & 3.215 & 3.252 \\
\bottomrule
\end{tabular}
\end{table}

RSE's 95\% CI ($[3.103, 3.115]$) is entirely non-overlapping with all other models, providing unambiguous evidence of superiority on simulated data.

\subsection{Key Findings}

\begin{enumerate}[noitemsep,leftmargin=*]
    \item \textbf{RSE achieves best tail risk:} CVaR$_{95} = 3.109$, $-3.29\%$ vs.\ LSTM ($p < 0.0001$, $d = -7.41$). Lowest CV (0.30\%) confirms ensemble stability.
    \item \textbf{Baselines show excellent stability:} LSTM, Transformer, 3SCH, W-DRO-T all have CV $< 1\%$, confirming the two-stage protocol's robustness.
    \item \textbf{DRO does not degrade in-distribution:} W-DRO-T CVaR $3.227$ vs Transformer $3.234$ ($p = 0.188$). The true benefit emerges under distributional shift.
    \item \textbf{3SCH preserves baseline performance:} CVaR $3.219$ vs LSTM $3.215$ ($p = 0.438$). The mixed-loss transition does not degrade performance.
    \item \textbf{RVSN and SAC-CVaR anomalous:} RVSN (CV = 60\%, positive P\&L) and SAC-CVaR (CVaR = 14.9) require algorithmic fixes.
\end{enumerate}

%===============================================================
\section{Real Market Validation and Crisis Testing}
\label{sec:crisis}

\subsection{Market Calibration Procedure}

Models are evaluated on market-calibrated Heston simulations (5,000 paths, 30-day horizon). Heston parameters calibrated to historical implied volatility surfaces from yfinance data. Two markets tested:
\begin{itemize}[noitemsep]
    \item \textbf{SPY} (3~bps cost, 2~bps slippage): Normal ($\sigma \approx 15\%$), Post-COVID ($20\%$), COVID ($65\%$), 2008 ($80\%$)
    \item \textbf{NIFTY} (18~bps cost, 5~bps slippage): Normal ($14\%$), COVID ($55\%$), 2008 ($70\%$)
\end{itemize}

\subsection{US Market Experiments (SPY)}

\begin{table}[H]
\centering\small
\caption{CVaR$_{95}$ across SPY scenarios (5,000 paths, 3 bps cost).}
\label{tab:spy}
\begin{tabular}{@{}lrrrr@{}}
\toprule
\textbf{Model} & \textbf{Normal} & \textbf{Post-COVID} & \textbf{COVID} & \textbf{GFC 2008} \\\midrule
LSTM & 20.28 & 27.36 & 86.84 & 100.88 \\
Transformer & 14.47 & 19.66 & 64.80 & 72.71 \\
RSE & 16.90 & 23.20 & 69.76 & 84.57 \\
\textbf{W-DRO-T} & \textbf{11.98} & \textbf{20.69} & \textbf{46.94} & \textbf{62.26} \\
3SCH & 15.79 & 21.43 & 69.25 & 88.43 \\
\bottomrule
\end{tabular}
\end{table}

W-DRO-T dominates all US scenarios: 41\% improvement over LSTM in normal, 46\% in COVID, 38\% in 2008. The DRO gradient penalty provides measurable distributional robustness precisely when capital is most constrained.

\subsection{Indian Market Experiments (NIFTY)}

\begin{table}[H]
\centering\small
\caption{CVaR$_{95}$ across NIFTY scenarios (5,000 paths, 18 bps cost).}
\label{tab:nifty}
\begin{tabular}{@{}lrrr@{}}
\toprule
\textbf{Model} & \textbf{Normal} & \textbf{COVID} & \textbf{GFC 2008} \\\midrule
LSTM & 785.62 & 3028.93 & 3645.62 \\
Transformer & 769.61 & \textbf{1644.19} & 2885.31 \\
RSE & 686.62 & 2592.97 & 2841.25 \\
W-DRO-T & 663.92 & 2883.96 & \textbf{2349.72} \\
\textbf{3SCH} & \textbf{622.26} & 2465.71 & 2962.35 \\
\bottomrule
\end{tabular}
\end{table}

Market-specific patterns: 3SCH achieves best NIFTY normal (622.26) due to low trading volume (0.90) minimizing 18~bps costs. Transformer dominates NIFTY COVID (1644.19). W-DRO-T best in 2008 stress (2349.72). The higher NIFTY costs penalize W-DRO-T's active trading.

3SCH's advantage stems from trading volume: 0.90 vs W-DRO-T 2.70 and Transformer 2.09. At 18 bps, each trade costs $6\times$ more than SPY:
\begin{itemize}[noitemsep]
    \item vs W-DRO-T (663.92): $-6.3\%$
    \item vs RSE (686.62): $-9.4\%$
    \item vs Transformer (769.61): $-19.1\%$
    \item vs LSTM (785.62): $-20.8\%$
\end{itemize}

\subsection{Crisis Robustness Analysis}

\begin{table}[H]
\centering\small
\caption{Crisis-to-normal CVaR$_{95}$ ratio (SPY; lower = more robust).}
\label{tab:robustness}
\begin{tabular}{@{}lccl@{}}
\toprule
\textbf{Model} & \textbf{COVID/Normal} & \textbf{2008/Normal} & \textbf{Interpretation} \\\midrule
\textbf{W-DRO-T} & $\mathbf{3.92\times}$ & $5.20\times$ & Best COVID robustness \\
RSE & $4.13\times$ & $5.00\times$ & Moderate \\
LSTM & $4.28\times$ & $4.97\times$ & Baseline \\
3SCH & $4.39\times$ & $5.60\times$ & Moderate \\
Transformer & $4.48\times$ & $5.02\times$ & High degradation \\
\bottomrule
\end{tabular}
\end{table}

W-DRO-T achieves the lowest COVID-to-normal ratio ($3.92\times$ vs LSTM $4.28\times$), confirming DRO regularization provides measurable distributional robustness under volatility regime shifts.

\subsection{P\&L Volatility Analysis}

\begin{table}[H]
\centering\small
\caption{Std P\&L across SPY scenarios.}
\label{tab:std_pnl}
\begin{tabular}{@{}lrrrr@{}}
\toprule
\textbf{Model} & \textbf{Normal} & \textbf{COVID} & \textbf{2008} & \textbf{Ratio} \\\midrule
\textbf{W-DRO-T} & \textbf{1.71} & \textbf{6.22} & \textbf{11.76} & $3.64\times$ \\
Transformer & 3.13 & 21.29 & 14.65 & $6.80\times$ \\
RSE & 3.57 & 15.03 & 19.47 & $4.21\times$ \\
3SCH & 3.70 & 16.38 & 21.06 & $4.43\times$ \\
LSTM & 5.05 & 22.18 & 25.88 & $4.39\times$ \\
\bottomrule
\end{tabular}
\end{table}

W-DRO-T achieves the lowest P\&L volatility across all scenarios: Std = 1.71 in normal (vs 5.05 for LSTM, $-66\%$) and 6.22 in COVID (vs 22.18, $-72\%$). The crisis-to-normal Std ratio of $3.64\times$ is also best, confirming DRO regularization provides smoother P\&L distributions.
